\documentclass[11pt,a4paper]{article}

\usepackage[utf8]{inputenc}
\usepackage[T1]{fontenc}
\usepackage[a4paper,margin=2.5cm]{geometry}
\usepackage[table]{xcolor}
\usepackage{titlesec}
\usepackage{fancyhdr}
\usepackage{enumitem}
\usepackage{longtable}
\usepackage{booktabs}
\usepackage{graphicx}
\usepackage{float}
\usepackage{hyperref}
\usepackage{parskip}
\usepackage{helvet}
\renewcommand{\familydefault}{\sfdefault}

\definecolor{accent}{HTML}{1F4E78}

\titleformat{\section}{\color{accent}\Large\bfseries}{\thesection.}{0.5em}{}
\titleformat{\subsection}{\color{accent}\large\bfseries}{\thesubsection}{0.5em}{}
\titlespacing*{\section}{0pt}{18pt}{8pt}

\pagestyle{fancy}
\fancyhf{}
\renewcommand{\headrulewidth}{0.4pt}
\fancyhead[L]{\small Compos\'es polyazot\'es -- bibliographie}
\fancyhead[R]{\small\thepage}

\hypersetup{colorlinks=true, linkcolor=accent, urlcolor=accent, citecolor=accent}

\title{\color{accent}\textbf{Compos\'es polyazot\'es (N$_x$) : neutres, cations et anions\\
\large Recueil bibliographique des structures et \'energies}}
\author{Gilles Frapper \\ Professeur des Universit\'es en Chimie Th\'eorique \\ IC2MP UMR 7285 CNRS, Universit\'e de Poitiers \\[2pt] \small \href{mailto:gilles.frapper@univ-poitiers.fr}{gilles.frapper@univ-poitiers.fr}}
\date{\today}

\begin{document}
\maketitle
\thispagestyle{fancy}

\begin{abstract}
\noindent Ce document recense \textbf{109 compos\'es polyazot\'es}
(N$_x$, $x=3$ \`a 20, exclusivement des clusters N -- aucun autre
\'el\'ement) rapport\'es dans \textbf{21 articles} de la litt\'erature,
organis\'es par \'etat de charge (neutre, cation, anion). Pour chaque
compos\'e : nom, groupe ponctuel de sym\'etrie, \'energie \'electronique
GFN2-xTB (\'eV/atome, niveau de th\'eorie commun \`a tout le recueil,
recalcul\'e ici pour permettre la comparaison malgr\'e l'h\'et\'erog\'en\'eit\'e
des niveaux de calcul originaux -- CCSD(T), B3LYP, MP2, HF selon
l'article), \'energie relative au sein du groupe (m\^eme formule, m\^eme
charge) en kcal/mol, et r\'ef\'erence bibliographique. Toutes les
g\'eom\'etries ont \'et\'e v\'erifi\'ees comme \'etant de vrais minima
locaux (analyse de fr\'equences, aucun mode imaginaire significatif).
\end{abstract}

\section{M\'ethode}

Chaque g\'eom\'etrie a \'et\'e reconstruite \`a partir des param\`etres
publi\'es (longueurs de liaison, angles, ou coordonn\'ees cart\'esiennes
selon ce que rapporte chaque article), puis relax\'ee en GFN2-xTB (via le
binding Python \texttt{tblite}, charge et multiplicit\'e de spin impos\'ees
telles que rapport\'ees) et v\'erifi\'ee par calcul de fr\'equences
(Hessienne num\'erique ; tout mode imaginaire significatif est suivi et la
structure r\'eoptimis\'ee, jusqu'\`a confirmation d'un vrai minimum). Ce
niveau de th\'eorie commun permet de comparer directement des structures
initialement rapport\'ees \`a des niveaux de calcul tr\`es diff\'erents --
un point essentiel puisque les articles sources vont de HF/6-31G* (1990s)
\`a CCSD(T)/aug-cc-pVTZ.

L'\'energie relative $E_{\mathrm{rel}}$ est calcul\'ee, pour chaque groupe
(m\^eme formule N$_x$, m\^eme charge), par rapport \`a l'isom\`ere le plus
stable du groupe ($E_{\mathrm{rel}} = 0$). Elle n'est donc compos\'ee que
\textbf{dans un m\^eme groupe} -- une comparaison entre formules ou entre
\'etats de charge diff\'erents n'est pas physiquement significative \`a ce
niveau (voir la discussion sur les r\'ef\'erences de formation dans le
rapport m\'ethodologique compagnon).

Sources : les 69 structures purement azot\'ees des r\'ef\'erences [GS] et
[B] (corpus biblio\_polyN initial, 70 structures moins le pentazole
N$_5$H qui contient de l'hydrog\`ene), \'etendues par 34 structures
suppl\'ementaires
issues de 19 articles ind\'ependants ([A1] \`a [A19]) -- parmi lesquelles 3
sont de nouvelles topologies absentes du corpus initial, les 31 autres
\'etant des confirmations ind\'ependantes (\`a un niveau de calcul
diff\'erent, parfois exp\'erimental) de topologies d\'ej\`a document\'ees.

\section{Compos\'es neutres}

\begin{longtable}{@{}p{4.4cm}lp{1.5cm}rrc@{}}
\caption{Compos\'es polyazot\'es \textbf{neutres} document\'es dans la bibliographie : nom, groupe ponctuel, \'energie GFN2-xTB (\'eV/atome), \'energie relative au sein du groupe (formule) en kcal/mol, r\'ef\'erence.}
\label{tab:final-neutre}\\
\toprule
\textbf{Nom} & \textbf{Formule} & \textbf{Sym.} & \textbf{$E$ (eV/at.)} & \textbf{$E_{\mathrm{rel}}$ (kcal/mol)} & \textbf{R\'ef.} \\
\midrule
\endfirsthead
\multicolumn{6}{c}{\small (suite)}\\
\toprule
\textbf{Nom} & \textbf{Formule} & \textbf{Sym.} & \textbf{$E$ (eV/at.)} & \textbf{$E_{\mathrm{rel}}$ (kcal/mol)} & \textbf{R\'ef.} \\
\midrule
\endhead
\bottomrule
\endfoot
\bottomrule
\endlastfoot
\addlinespace
\texttt{N2} & N$_{2}$ & Dinfh & -78{,}4223 & 0{,}00 & [B] \\
\addlinespace
\texttt{N3\_radical} & N$_{3}$ & Dinfh & -78{,}2386 & 0{,}00 & [B] \\
\addlinespace
\texttt{N4\_C2v\_butterfly} & N$_{4}$ & C2v & -78{,}2786 & 0{,}00 & [B] \\
\texttt{N4\_C2h\_chain} & N$_{4}$ & C2h & -78{,}2786 & 0{,}00 & [B] \\
\texttt{N4\_Cs\_triplet\_4\_openchain} & N$_{4}$ & Cs & -78{,}2786 & 0{,}00 & [A4] \\
\texttt{N4\_Cs\_triplet\_5\_openchain} & N$_{4}$ & Cs & -78{,}2786 & 0{,}00 & [A6] \\
\texttt{N4\_C2h\_chain\_GS} & N$_{4}$ & C2h & -78{,}2786 & 0{,}00 & [GS] \\
\texttt{N4\_Td} & N$_{4}$ & Td & -78{,}1261 & 14{,}06 & [B] \\
\texttt{N4\_Td\_tetrahedrane\_CCSDT} & N$_{4}$ & Td & -78{,}1261 & 14{,}06 & [A3] \\
\texttt{N4\_Td\_singlet\_1} & N$_{4}$ & Td & -78{,}1261 & 14{,}06 & [A4] \\
\texttt{N4\_Td\_tetraazatetrahedrane} & N$_{4}$ & Td & -78{,}1261 & 14{,}06 & [A6] \\
\texttt{N4\_Td\_GS} & N$_{4}$ & Td & -78{,}1261 & 14{,}06 & [GS] \\
\texttt{N4\_D2h\_rectangle\_planar} & N$_{4}$ & D2h & -77{,}6748 & 55{,}70 & [B] \\
\texttt{N4\_D2h\_rectangle\_cyclobutadiene\_analog} & N$_{4}$ & D2h & -77{,}6748 & 55{,}70 & [A5] \\
\texttt{N4\_D2h\_tetrazete\_GS} & N$_{4}$ & D2h & -77{,}6748 & 55{,}70 & [GS] \\
\texttt{N4\_D2d\_puckered} & N$_{4}$ & D2d & -77{,}3153 & 88{,}86 & [B] \\
\texttt{N4\_D2d\_triplet\_5} & N$_{4}$ & D2d & -77{,}3153 & 88{,}86 & [A4] \\
\addlinespace
\texttt{N5\_C2v\_radical\_2A1} & N$_{5}$ & C2v & -78{,}6104 & 0{,}00 & [A7] \\
\texttt{N5\_C2v\_radical\_2B2} & N$_{5}$ & C2v & -78{,}6104 & 0{,}00 & [A7] \\
\addlinespace
\texttt{N6\_C2h\_chain} & N$_{6}$ & C2h & -78{,}8828 & 0{,}00 & [B] \\
\texttt{N6\_C2h\_trans\_diazide\_v2} & N$_{6}$ & C2h & -78{,}8828 & 0{,}00 & [A9] \\
\texttt{N6\_C2h\_diazide\_synthesis\_product} & N$_{6}$ & C2h & -78{,}8828 & 0{,}00 & [A16] \\
\texttt{N6\_C2h\_openchain\_GS} & N$_{6}$ & C2h & -78{,}8828 & 0{,}00 & [GS] \\
\texttt{N6\_D2\_twistboat\_GS} & N$_{6}$ & D2 & -78{,}7087 & 24{,}08 & [GS] \\
\texttt{N6\_C2\_book} & N$_{6}$ & C2 & -78{,}7087 & 24{,}09 & [B] \\
\texttt{N6\_D2\_twisted} & N$_{6}$ & D2 & -78{,}7087 & 24{,}09 & [B] \\
\texttt{N6\_D6h\_hexagon\_aromatic} & N$_{6}$ & D6h & -78{,}7087 & 24{,}09 & [A5] \\
\texttt{N6\_D3h\_prism} & N$_{6}$ & D3h & -77{,}8951 & 136{,}66 & [B] \\
\addlinespace
\texttt{N7\_C2\_linearZZ} & N$_{7}$ & C2 & -78{,}8040 & 0{,}00 & [B] \\
\texttt{N7\_iso7\_Cs\_chain} & N$_{7}$ & Cs & -78{,}8040 & 0{,}00 & [A10] \\
\texttt{N7\_iso7\_C2\_chain} & N$_{7}$ & C2 & -78{,}8040 & 0{,}00 & [A10] \\
\texttt{N7\_wang\_3\_C2v\_openchain\_DUPLICATE} & N$_{7}$ & C2v & -78{,}8040 & 0{,}00 & [A17] \\
\texttt{N7\_Cs\_ring\_chain} & N$_{7}$ & Cs & -78{,}6998 & 16{,}83 & [B] \\
\texttt{N7\_C2v\_logcarrier} & N$_{7}$ & C2v & -78{,}5721 & 37{,}45 & [B] \\
\addlinespace
\texttt{N8\_C2h\_ladder} & N$_{8}$ & C2h & -79{,}0369 & 0{,}00 & [B] \\
\texttt{N8\_C2v\_ring} & N$_{8}$ & C2v & -79{,}0369 & 0{,}00 & [B] \\
\texttt{N8\_D2h\_pentalene} & N$_{8}$ & D2h & -79{,}0369 & 0{,}00 & [GS] \\
\texttt{N8\_D2d\_octaazacyclooctatetraene} & N$_{8}$ & D2d & -79{,}0369 & 0{,}00 & [GS] \\
\texttt{N8\_pentalene\_dianion\_analog\_struct11} & N$_{8}$ & n.d. & -79{,}0369 & 0{,}00 & [A5] \\
\texttt{N8\_D2h\_octaazapentalene} & N$_{8}$ & D2h & -79{,}0369 & 0{,}00 & [A14] \\
\texttt{N8\_D2h\_pentalene\_v3} & N$_{8}$ & D2h & -79{,}0369 & 0{,}00 & [A15] \\
\texttt{N8\_Cs\_pentagonal} & N$_{8}$ & Cs & -78{,}9832 & 9{,}92 & [B] \\
\texttt{N8\_Cs\_azidopentazole} & N$_{8}$ & Cs & -78{,}9832 & 9{,}92 & [GS] \\
\texttt{N8\_C2h\_ZEZ} & N$_{8}$ & C2h & -78{,}9271 & 20{,}26 & [B] \\
\texttt{N8\_C2v\_EZE} & N$_{8}$ & C2v & -78{,}9271 & 20{,}26 & [B] \\
\texttt{N8\_D4h\_cyclooctatetraene} & N$_{8}$ & n.d. & -78{,}8846 & 28{,}10 & [A5] \\
\texttt{N8\_C1\_ZZE} & N$_{8}$ & C1 & -78{,}8775 & 29{,}41 & [B] \\
\texttt{N8\_Cs\_ZEE} & N$_{8}$ & Cs & -78{,}8775 & 29{,}41 & [B] \\
\texttt{N8\_C2h\_EEE} & N$_{8}$ & C2h & -78{,}8775 & 29{,}41 & [B] \\
\texttt{N8\_LiWang\_C2h\_openchain} & N$_{8}$ & C2h & -78{,}8775 & 29{,}41 & [A13] \\
\texttt{N8\_Cs\_azidopentazole\_v2} & N$_{8}$ & Cs & -78{,}8775 & 29{,}41 & [A14] \\
\texttt{N8\_C2h\_diazidodiimide\_cis\_3a} & N$_{8}$ & C2h & -78{,}8775 & 29{,}41 & [A14] \\
\texttt{N8\_D2d\_cis\_cyclooctatetraene} & N$_{8}$ & D2d & -78{,}8775 & 29{,}41 & [A15] \\
\texttt{N8\_D2h\_ring\_pendant} & N$_{8}$ & D2h & -78{,}7585 & 51{,}36 & [B] \\
\texttt{N8\_C2v\_branched} & N$_{8}$ & C2v & -78{,}3470 & 127{,}27 & [B] \\
\texttt{N8\_Oh\_cube} & N$_{8}$ & Oh & -77{,}9219 & 205{,}70 & [B] \\
\texttt{N8\_Oh\_cubane\_GS} & N$_{8}$ & Oh & -77{,}9219 & 205{,}70 & [GS] \\
\texttt{N8\_Oh\_octaazacubane\_CCSD} & N$_{8}$ & Oh & -77{,}9219 & 205{,}70 & [A3] \\
\texttt{N8\_Oh\_cubane\_v2} & N$_{8}$ & Oh & -77{,}9219 & 205{,}70 & [A15] \\
\addlinespace
\texttt{N9\_chain} & N$_{9}$ & C2v & -78{,}8426 & 0{,}00 & [B] \\
\texttt{N9\_fused\_rings} & N$_{9}$ & C2v & -78{,}8324 & 2{,}13 & [B] \\
\addlinespace
\texttt{N10\_D2d\_linked5rings} & N$_{10}$ & D2d & -79{,}0043 & 0{,}00 & [B] \\
\texttt{N10\_D2d\_bispentazole\_GS} & N$_{10}$ & D2d & -79{,}0043 & 0{,}00 & [GS] \\
\texttt{N10\_cage\_C2v} & N$_{10}$ & C2v & -78{,}8978 & 24{,}56 & [B] \\
\texttt{N10\_C3\_cap} & N$_{10}$ & C3 & -78{,}8978 & 24{,}56 & [B] \\
\texttt{N10\_Cs\_ring\_chain} & N$_{10}$ & Cs & -78{,}8593 & 33{,}43 & [B] \\
\texttt{N10\_D10h\_ring} & N$_{10}$ & D10h & -78{,}8369 & 38{,}61 & [GS] \\
\texttt{N10\_D5h\_prism} & N$_{10}$ & D5h & -78{,}2066 & 183{,}95 & [B] \\
\addlinespace
\texttt{N11\_C2v\_acyclic\_chain}$^{\dagger}$ & N$_{11}$ & C2v & -78{,}8583 & 0{,}00 & [A18] \\
\addlinespace
\texttt{N12\_C2h\_dipentazolyldiazene} & N$_{12}$ & C2h & -78{,}3628 & 0{,}00 & [GS] \\
\addlinespace
\texttt{N18\_C2v\_cage} & N$_{18}$ & C2v & -78{,}6312 & 0{,}00 & [A19] \\
\addlinespace
\texttt{N20\_Ih\_dodecahedrane} & N$_{20}$ & Ih & -78{,}6645 & 0{,}00 & [GS] \\
\texttt{N20\_Ih\_cage\_HaEtAl} & N$_{20}$ & Ih & -78{,}6645 & 0{,}00 & [A12] \\
\multicolumn{6}{@{}p{15.5cm}}{\footnotesize $^{\dagger}$ Reconstruction de confiance basse (repli g\'en\'erique faute de connectivit\'e compl\`ete dans l'article source) -- minimum GFN2-xTB v\'erifi\'e, mais rien ne garantit la co\"incidence avec l'isom\`ere pr\'ecis d\'ecrit dans l'article.}\\
\end{longtable}


\section{Compos\'es cationiques}

\begin{longtable}{@{}p{4.4cm}lp{1.5cm}rrc@{}}
\caption{Compos\'es polyazot\'es \textbf{cationiques} document\'es dans la bibliographie.}
\label{tab:final-cation}\\
\toprule
\textbf{Nom} & \textbf{Formule} & \textbf{Sym.} & \textbf{$E$ (eV/at.)} & \textbf{$E_{\mathrm{rel}}$ (kcal/mol)} & \textbf{R\'ef.} \\
\midrule
\endfirsthead
\multicolumn{6}{c}{\small (suite)}\\
\toprule
\textbf{Nom} & \textbf{Formule} & \textbf{Sym.} & \textbf{$E$ (eV/at.)} & \textbf{$E_{\mathrm{rel}}$ (kcal/mol)} & \textbf{R\'ef.} \\
\midrule
\endhead
\bottomrule
\endfoot
\bottomrule
\endlastfoot
\addlinespace
\texttt{N2+} & N$_{2}^{+}$ & Dinfh & -69{,}4370 & 0{,}00 & [B] \\
\addlinespace
\texttt{N3+\_linear} & N$_{3}^{+}$ & Dinfh & -72{,}5478 & 0{,}00 & [B] \\
\texttt{N3+\_Dinfh\_linear\_groundstate} & N$_{3}^{+}$ & Dinfh & -72{,}5239 & 1{,}65 & [A1] \\
\texttt{N3+\_triangle} & N$_{3}^{+}$ & D3h & -72{,}1692 & 26{,}19 & [B] \\
\addlinespace
\texttt{N4+\_rectangle} & N$_{4}^{+}$ & D2h & -73{,}6181 & 0{,}00 & [B] \\
\addlinespace
\texttt{N5+\_Vchain} & N$_{5}^{+}$ & C2v & -76{,}1117 & 0{,}00 & [B] \\
\texttt{N5+\_C2v\_experimental\_Xray} & N$_{5}^{+}$ & C2v & -76{,}1117 & 0{,}00 & [A8] \\
\addlinespace
\texttt{N6+\_C2h\_planar} & N$_{6}^{+}$ & C2h & -76{,}4957 & 0{,}00 & [B] \\
\texttt{N6+\_C2v} & N$_{6}^{+}$ & C2v & -76{,}4957 & 0{,}00 & [B] \\
\texttt{N6+\_C2h\_nonplanar} & N$_{6}^{+}$ & C2h & -76{,}4957 & 0{,}00 & [B] \\
\addlinespace
\texttt{N7+\_chain} & N$_{7}^{+}$ & C2v & -76{,}9440 & 0{,}00 & [B] \\
\texttt{N7+\_1\_C2v\_openchain} & N$_{7}^{+}$ & C2v & -76{,}9440 & 0{,}00 & [A11] \\
\addlinespace
\texttt{N9+\_chain} & N$_{9}^{+}$ & C2v & -77{,}5122 & 0{,}00 & [B] \\
\texttt{N9+\_fused\_rings} & N$_{9}^{+}$ & C2v & -77{,}3283 & 38{,}17 & [B] \\
\addlinespace
\texttt{N10+\_chain} & N$_{10}^{+}$ & C1 & -77{,}6183 & 0{,}00 & [B] \\
\texttt{N10+\_rings} & N$_{10}^{+}$ & C1 & -77{,}4938 & 28{,}71 & [B] \\
\end{longtable}


\section{Compos\'es anioniques}

\begin{longtable}{@{}p{4.4cm}lp{1.5cm}rrc@{}}
\caption{Compos\'es polyazot\'es \textbf{anioniques} document\'es dans la bibliographie.}
\label{tab:final-anion}\\
\toprule
\textbf{Nom} & \textbf{Formule} & \textbf{Sym.} & \textbf{$E$ (eV/at.)} & \textbf{$E_{\mathrm{rel}}$ (kcal/mol)} & \textbf{R\'ef.} \\
\midrule
\endfirsthead
\multicolumn{6}{c}{\small (suite)}\\
\toprule
\textbf{Nom} & \textbf{Formule} & \textbf{Sym.} & \textbf{$E$ (eV/at.)} & \textbf{$E_{\mathrm{rel}}$ (kcal/mol)} & \textbf{R\'ef.} \\
\midrule
\endhead
\bottomrule
\endfoot
\bottomrule
\endlastfoot
\addlinespace
\texttt{N3-\_linear} & N$_{3}^{-}$ & Dinfh & -80{,}6963 & 0{,}00 & [B] \\
\texttt{N3-\_bent} & N$_{3}^{-}$ & C2v & -79{,}6529 & 72{,}18 & [B] \\
\texttt{N3-\_triangle} & N$_{3}^{-}$ & D3h & -79{,}3892 & 90{,}43 & [B] \\
\addlinespace
\texttt{N4-\_rectangle} & N$_{4}^{-}$ & D2h & -79{,}6030 & 0{,}00 & [B] \\
\texttt{N4-\_chain\_planar} & N$_{4}^{-}$ & C2h & -79{,}2985 & 28{,}09 & [B] \\
\addlinespace
\texttt{N5-\_pentagon} & N$_{5}^{-}$ & D5h & -80{,}3329 & 0{,}00 & [B] \\
\texttt{N5-\_pentagon\_D5h\_crosscheck} & N$_{5}^{-}$ & D5h & -80{,}3329 & 0{,}00 & [A2] \\
\texttt{N5-\_D5h\_pentazolate} & N$_{5}^{-}$ & D5h & -80{,}3329 & 0{,}00 & [A7] \\
\texttt{N5-\_bipyramid} & N$_{5}^{-}$ & D3h & -79{,}0896 & 143{,}36 & [B] \\
\addlinespace
\texttt{N6-\_Cs\_chain} & N$_{6}^{-}$ & Cs & -79{,}8151 & 0{,}00 & [B] \\
\texttt{N6-\_C2\_linked\_triangles} & N$_{6}^{-}$ & C2 & -79{,}0561 & 105{,}02 & [B] \\
\addlinespace
\texttt{N7-\_chain} & N$_{7}^{-}$ & C2v & -79{,}7741 & 0{,}00 & [B] \\
\texttt{N7-\_6\_Cs\_6ring\_boat}$^{\dagger}$ & N$_{7}^{-}$ & Cs & -79{,}7176 & 9{,}12 & [A11] \\
\addlinespace
\texttt{N8-\_C2h\_chain} & N$_{8}^{-}$ & C2h & -79{,}6681 & 0{,}00 & [B] \\
\texttt{N8-\_C2h\_ZEZ} & N$_{8}^{-}$ & C2h & -79{,}6681 & 0{,}00 & [B] \\
\addlinespace
\texttt{N9-\_chain\_twistedW} & N$_{9}^{-}$ & C2 & -79{,}6546 & 0{,}00 & [B] \\
\texttt{N9-\_chain} & N$_{9}^{-}$ & C2v & -79{,}6546 & 0{,}00 & [B] \\
\texttt{N9-\_ring\_chain} & N$_{9}^{-}$ & Cs & -79{,}6416 & 2{,}69 & [B] \\
\addlinespace
\texttt{N10-\_D2h\_perp\_rings} & N$_{10}^{-}$ & D2h & -79{,}7228 & 0{,}00 & [B] \\
\multicolumn{6}{@{}p{15.5cm}}{\footnotesize $^{\dagger}$ Reconstruction de confiance basse (repli g\'en\'erique faute de connectivit\'e compl\`ete dans l'article source) -- minimum GFN2-xTB v\'erifi\'e, mais rien ne garantit la co\"incidence avec l'isom\`ere pr\'ecis d\'ecrit dans l'article.}\\
\end{longtable}


\section{Limites m\'ethodologiques : fiabilit\'e de l'\'energie relative GFN2-xTB}
\label{sec:limites-gfn2}

Le niveau GFN2-xTB, utilis\'e syst\'ematiquement dans ce recueil pour
disposer d'une \'echelle d'\'energie commune \`a des structures publi\'ees
\`a des niveaux de calcul tr\`es h\'et\'erog\`enes, a \'et\'e confront\'e
directement aux propres valeurs de r\'ef\'erence de [B] (tableaux
``Heats of Formation'', p.73-75, niveaux CCSD(T)/G2/B3LYP selon les cas)
pour les formules o\`u plusieurs isom\`eres sont document\'es des deux
c\^ot\'es. Le constat est net : \textbf{GFN2-xTB reproduit correctement les
grands \'ecarts (structures tr\`es instables comme le cubane N$_8$
O$_h$) mais se r\'ev\`ele peu fiable, et parfois qualitativement faux, sur
les \'ecarts fins ou moyens (10 \`a 70 kcal/mol) entre isom\`eres proches}
:

\begin{table}[H]
\centering
\small
\begin{tabular}{@{}p{1.3cm}p{5.4cm}rr@{}}
\toprule
Formule & Comparaison & \'Ecart [B] (kcal/mol) & \'Ecart GFN2-xTB (kcal/mol) \\
\midrule
N$_4$ & chaine C2h ($^3B_u$) vs papillon C2v & 73{,}0 & \textbf{0{,}0 (effondr\'es)} \\
N$_6$ & ``book'' C2 vs ``twisted'' D2 & 36{,}3 & \textbf{0{,}0 (effondr\'es)} \\
N$_7$ & 5-ring+chain (Cs) vs cha\^ine lin\'eaire ZZ (C2) & 27{,}0 (Cs plus stable) & 16{,}8 (\textbf{ordre invers\'e}) \\
N$_8$ & ``ladder'' C2h : rang parmi 11 isom\`eres & avant-dernier (317{,}0/422{,}7) & \textbf{premier (ex \ae quo, 0{,}0)} \\
N$_3^{+}$ & triangle D3h vs lin\'eaire D$_{\infty h}$ & 9{,}6 (triangle plus stable) & 26{,}2 (\textbf{ordre invers\'e}) \\
N$_5^{-}$ & pentagone D5h vs bipyramide D3h & 189{,}6 & 143{,}4 (sous-estim\'e de 46) \\
\bottomrule
\end{tabular}
\caption{Comparaison directe des \'ecarts d'\'energie relative entre
isom\`eres de m\^eme formule/charge, GFN2-xTB (ce recueil) contre les
valeurs propres de [B] au niveau CCSD(T)/DFT.}
\label{tab:gfn2-vs-bartlett}
\end{table}

Cons\'equence pratique : la colonne $E_{\mathrm{rel}}$ des tableaux 1 \`a
3 (compos\'es neutres/cationiques/anioniques) et les $\Delta H_f$ de la
section~\ref{sec:isodesmic} qui suit doivent \^etre lus comme un
\textbf{classement de d\'egrossissage} (fiable pour distinguer un
isom\`ere tr\`es instable d'un isom\`ere comp\'etitif) et non comme une
\'energie relative quantitativement fiable entre isom\`eres proches --
deux isom\`eres donn\'es comme ex \ae quo ou proches dans ce recueil
peuvent en r\'ealit\'e diff\'erer de plusieurs dizaines de kcal/mol, et
l'ordre de stabilit\'e entre deux isom\`eres proches peut \^etre invers\'e
par rapport \`a un calcul de plus haut niveau. Une campagne de
raffinement DFT (wB97X-D3/def2-TZVP ou \'equivalent) puis
DLPNO-CCSD(T)/def2-TZVP, d\'ej\`a configur\'ee dans le d\'epot compagnon
\texttt{polyN\_adapt/refinement/} mais destin\'ee \`a s'ex\'ecuter sur un
calculateur distant, doit \^etre men\'ee avant d'utiliser ce recueil pour
trancher entre isom\`eres d'\'energie proche.

\section{Enthalpies de r\'eaction quasi-isodesmique}
\label{sec:isodesmic}

\`A la suite de [B] (p.75), l'enthalpie de formation de chaque compos\'e
est estim\'ee par une r\'eaction quasi-isodesmique r\'ef\'erenc\'ee \`a
l'azote mol\'eculaire (esp\`eces neutres) ou \`a un ion azot\'e-hydrog\'en\'e
de $\Delta H_f$ exp\'erimental connu (esp\`eces charg\'ees, afin
d'\'equilibrer la charge sans r\'ef\'erence \`a un \'electron libre) :

\begin{itemize}
  \item Neutres : N$_x \to$ (x/2)~N$_2$
  \item Cations : [NH$_4$]$^+$ + (x$-$1)/2~N$_2 \to$ N$_x^+$ + 2~H$_2$,
        avec $\Delta H_f$(NH$_4^+$) = 150{,}6 kcal/mol (exp\'erimental)
  \item Anions : [NH$_2$]$^-$ + (x$-$1)/2~N$_2 \to$ N$_x^-$ + H$_2$,
        avec $\Delta H_f$(NH$_2^-$) = 27{,}0 kcal/mol (exp\'erimental)
\end{itemize}

Toutes les \'energies \'electroniques (N$_x$, N$_2$, NH$_4^+$, H$_2$,
NH$_2^-$) sont calcul\'ees au m\^eme niveau GFN2-xTB que le reste de ce
recueil, avec les limites discut\'ees en section~\ref{sec:limites-gfn2} --
en particulier, les r\'eactions cationique/anionique cassent des liaisons
N--H et forment H--H, un type d'\'echange de liaison sur lequel GFN2-xTB
n'a pas \'et\'e valid\'e ici : les $\Delta H_f$ absolus obtenus pour les
ions peuvent s'\'ecarter des valeurs de [B] de plus de 100 kcal/mol (ex. :
N$_5^{+}$ calcul\'e ici \`a 195{,}9 kcal/mol contre 321{,}8 chez [B]). Ces
valeurs sont n\'eanmoins conserv\'ees ci-dessous, coh\'erentes en interne
avec le reste du recueil, en attendant leur raffinement par la campagne
DFT/CCSD(T) \'evoqu\'ee ci-dessus.

\subsection{Neutres}
\begin{longtable}{@{}p{4.6cm}lp{1.4cm}rc@{}}
\caption{Enthalpie de r\'eaction quasi-isodesmique (GFN2-xTB) des compos\'es \textbf{neutres}, r\'eaction N$_x \to$ (x/2)~N$_2$.}
\label{tab:isodesmic-neutre}\\
\toprule
\textbf{Nom} & \textbf{Formule} & \textbf{Sym.} & \textbf{$\Delta H_f$ (kcal/mol)} & \textbf{R\'ef.} \\
\midrule
\endfirsthead
\multicolumn{5}{c}{\small (suite)}\\
\toprule
\textbf{Nom} & \textbf{Formule} & \textbf{Sym.} & \textbf{$\Delta H_f$ (kcal/mol)} & \textbf{R\'ef.} \\
\midrule
\endhead
\bottomrule
\endfoot
\bottomrule
\endlastfoot
\addlinespace
\texttt{N2} & N$_{2}$ & Dinfh & -0{,}0 & [B] \\
\addlinespace
\texttt{N3\_radical} & N$_{3}$ & Dinfh & 12{,}7 & [B] \\
\addlinespace
\texttt{N4\_C2v\_butterfly} & N$_{4}$ & C2v & 13{,}3 & [B] \\
\texttt{N4\_C2h\_chain} & N$_{4}$ & C2h & 13{,}3 & [B] \\
\texttt{N4\_Cs\_triplet\_4\_openchain} & N$_{4}$ & Cs & 13{,}3 & [A4] \\
\texttt{N4\_Cs\_triplet\_5\_openchain} & N$_{4}$ & Cs & 13{,}3 & [A6] \\
\texttt{N4\_C2h\_chain\_GS} & N$_{4}$ & C2h & 13{,}3 & [GS] \\
\texttt{N4\_Td} & N$_{4}$ & Td & 27{,}3 & [B] \\
\texttt{N4\_Td\_tetrahedrane\_CCSDT} & N$_{4}$ & Td & 27{,}3 & [A3] \\
\texttt{N4\_Td\_singlet\_1} & N$_{4}$ & Td & 27{,}3 & [A4] \\
\texttt{N4\_Td\_tetraazatetrahedrane} & N$_{4}$ & Td & 27{,}3 & [A6] \\
\texttt{N4\_Td\_GS} & N$_{4}$ & Td & 27{,}3 & [GS] \\
\texttt{N4\_D2h\_rectangle\_planar} & N$_{4}$ & D2h & 69{,}0 & [B] \\
\texttt{N4\_D2h\_rectangle\_cyclobutadiene\_analog} & N$_{4}$ & D2h & 69{,}0 & [A5] \\
\texttt{N4\_D2h\_tetrazete\_GS} & N$_{4}$ & D2h & 69{,}0 & [GS] \\
\texttt{N4\_D2d\_puckered} & N$_{4}$ & D2d & 102{,}1 & [B] \\
\texttt{N4\_D2d\_triplet\_5} & N$_{4}$ & D2d & 102{,}1 & [A4] \\
\addlinespace
\texttt{N5\_C2v\_radical\_2A1} & N$_{5}$ & C2v & -21{,}7 & [A7] \\
\texttt{N5\_C2v\_radical\_2B2} & N$_{5}$ & C2v & -21{,}7 & [A7] \\
\addlinespace
\texttt{N6\_C2h\_chain} & N$_{6}$ & C2h & -63{,}7 & [B] \\
\texttt{N6\_C2h\_trans\_diazide\_v2} & N$_{6}$ & C2h & -63{,}7 & [A9] \\
\texttt{N6\_C2h\_diazide\_synthesis\_product} & N$_{6}$ & C2h & -63{,}7 & [A16] \\
\texttt{N6\_C2h\_openchain\_GS} & N$_{6}$ & C2h & -63{,}7 & [GS] \\
\texttt{N6\_D2\_twistboat\_GS} & N$_{6}$ & D2 & -39{,}6 & [GS] \\
\texttt{N6\_C2\_book} & N$_{6}$ & C2 & -39{,}6 & [B] \\
\texttt{N6\_D2\_twisted} & N$_{6}$ & D2 & -39{,}6 & [B] \\
\texttt{N6\_D6h\_hexagon\_aromatic} & N$_{6}$ & D6h & -39{,}6 & [A5] \\
\texttt{N6\_D3h\_prism} & N$_{6}$ & D3h & 73{,}0 & [B] \\
\addlinespace
\texttt{N7\_C2\_linearZZ} & N$_{7}$ & C2 & -61{,}6 & [B] \\
\texttt{N7\_iso7\_Cs\_chain} & N$_{7}$ & Cs & -61{,}6 & [A10] \\
\texttt{N7\_iso7\_C2\_chain} & N$_{7}$ & C2 & -61{,}6 & [A10] \\
\texttt{N7\_wang\_3\_C2v\_openchain\_DUPLICATE} & N$_{7}$ & C2v & -61{,}6 & [A17] \\
\texttt{N7\_Cs\_ring\_chain} & N$_{7}$ & Cs & -44{,}8 & [B] \\
\texttt{N7\_C2v\_logcarrier} & N$_{7}$ & C2v & -24{,}2 & [B] \\
\addlinespace
\texttt{N8\_C2h\_ladder} & N$_{8}$ & C2h & -113{,}4 & [B] \\
\texttt{N8\_C2v\_ring} & N$_{8}$ & C2v & -113{,}4 & [B] \\
\texttt{N8\_D2h\_pentalene} & N$_{8}$ & D2h & -113{,}4 & [GS] \\
\texttt{N8\_D2d\_octaazacyclooctatetraene} & N$_{8}$ & D2d & -113{,}4 & [GS] \\
\texttt{N8\_pentalene\_dianion\_analog\_struct11} & N$_{8}$ & n.d. & -113{,}4 & [A5] \\
\texttt{N8\_D2h\_octaazapentalene} & N$_{8}$ & D2h & -113{,}4 & [A14] \\
\texttt{N8\_D2h\_pentalene\_v3} & N$_{8}$ & D2h & -113{,}4 & [A15] \\
\texttt{N8\_D2h\_pentalene\_B} & N$_{8}$ & D2h & -113{,}4 & [B] \\
\texttt{N8\_Cs\_pentagonal} & N$_{8}$ & Cs & -103{,}5 & [B] \\
\texttt{N8\_Cs\_azidopentazole} & N$_{8}$ & Cs & -103{,}5 & [GS] \\
\texttt{N8\_C2h\_ZEZ} & N$_{8}$ & C2h & -93{,}1 & [B] \\
\texttt{N8\_C2v\_EZE} & N$_{8}$ & C2v & -93{,}1 & [B] \\
\texttt{N8\_D4h\_cyclooctatetraene} & N$_{8}$ & n.d. & -85{,}3 & [A5] \\
\texttt{N8\_C1\_ZZE} & N$_{8}$ & C1 & -84{,}0 & [B] \\
\texttt{N8\_Cs\_ZEE} & N$_{8}$ & Cs & -84{,}0 & [B] \\
\texttt{N8\_C2h\_EEE} & N$_{8}$ & C2h & -84{,}0 & [B] \\
\texttt{N8\_LiWang\_C2h\_openchain} & N$_{8}$ & C2h & -84{,}0 & [A13] \\
\texttt{N8\_Cs\_azidopentazole\_v2} & N$_{8}$ & Cs & -84{,}0 & [A14] \\
\texttt{N8\_C2h\_diazidodiimide\_cis\_3a} & N$_{8}$ & C2h & -84{,}0 & [A14] \\
\texttt{N8\_D2d\_cis\_cyclooctatetraene} & N$_{8}$ & D2d & -84{,}0 & [A15] \\
\texttt{N8\_D2h\_ring\_pendant} & N$_{8}$ & D2h & -62{,}0 & [B] \\
\texttt{N8\_C2v\_branched} & N$_{8}$ & C2v & 13{,}9 & [B] \\
\texttt{N8\_Oh\_cube} & N$_{8}$ & Oh & 92{,}3 & [B] \\
\texttt{N8\_Oh\_cubane\_GS} & N$_{8}$ & Oh & 92{,}3 & [GS] \\
\texttt{N8\_Oh\_octaazacubane\_CCSD} & N$_{8}$ & Oh & 92{,}3 & [A3] \\
\texttt{N8\_Oh\_cubane\_v2} & N$_{8}$ & Oh & 92{,}3 & [A15] \\
\addlinespace
\texttt{N9\_chain} & N$_{9}$ & C2v & -87{,}2 & [B] \\
\texttt{N9\_fused\_rings} & N$_{9}$ & C2v & -85{,}1 & [B] \\
\addlinespace
\texttt{N10\_D2d\_linked5rings} & N$_{10}$ & D2d & -134{,}2 & [B] \\
\texttt{N10\_D2d\_bispentazole\_GS} & N$_{10}$ & D2d & -134{,}2 & [GS] \\
\texttt{N10\_C1\_pentaazopentazole} & N$_{10}$ & C1 & -127{,}2 & [A20] \\
\texttt{N10\_cage\_C2v} & N$_{10}$ & C2v & -109{,}6 & [B] \\
\texttt{N10\_C3\_cap} & N$_{10}$ & C3 & -109{,}6 & [B] \\
\texttt{N10\_Cs\_ring\_chain} & N$_{10}$ & Cs & -100{,}8 & [B] \\
\texttt{N10\_D10h\_ring} & N$_{10}$ & D10h & -95{,}6 & [GS] \\
\texttt{N10\_D5h\_prism} & N$_{10}$ & D5h & 49{,}7 & [B] \\
\addlinespace
\texttt{N11\_C2v\_acyclic\_chain} & N$_{11}$ & C2v & -110{,}6 & [A18] \\
\addlinespace
\texttt{N12\_C2h\_dipentazolyldiazene} & N$_{12}$ & C2h & 16{,}5 & [GS] \\
\addlinespace
\texttt{N18\_C2v\_cage} & N$_{18}$ & C2v & -86{,}7 & [A19] \\
\addlinespace
\texttt{N20\_Ih\_dodecahedrane} & N$_{20}$ & Ih & -111{,}7 & [GS] \\
\texttt{N20\_Ih\_cage\_HaEtAl} & N$_{20}$ & Ih & -111{,}7 & [A12] \\
\end{longtable}


\subsection{Cations}
\begin{longtable}{@{}p{4.6cm}lp{1.4cm}rc@{}}
\caption{Enthalpie de r\'eaction quasi-isodesmique (GFN2-xTB) des compos\'es \textbf{cationiques}, r\'ef\'erenc\'ee \`a NH$_4^+$ ($\Delta H_f=150{,}6$ kcal/mol, valeur exp\'erimentale).}
\label{tab:isodesmic-cation}\\
\toprule
\textbf{Nom} & \textbf{Formule} & \textbf{Sym.} & \textbf{$\Delta H_f$ (kcal/mol)} & \textbf{R\'ef.} \\
\midrule
\endfirsthead
\multicolumn{5}{c}{\small (suite)}\\
\toprule
\textbf{Nom} & \textbf{Formule} & \textbf{Sym.} & \textbf{$\Delta H_f$ (kcal/mol)} & \textbf{R\'ef.} \\
\midrule
\endhead
\bottomrule
\endfoot
\bottomrule
\endlastfoot
\addlinespace
\texttt{N2+} & N$_{2}^{+}$ & Dinfh & 343{,}8 & [B] \\
\addlinespace
\texttt{N3+\_linear} & N$_{3}^{+}$ & Dinfh & 335{,}8 & [B] \\
\texttt{N3+\_Dinfh\_linear\_groundstate} & N$_{3}^{+}$ & Dinfh & 337{,}5 & [A1] \\
\texttt{N3+\_triangle} & N$_{3}^{+}$ & D3h & 362{,}0 & [B] \\
\addlinespace
\texttt{N4+\_rectangle} & N$_{4}^{+}$ & D2h & 372{,}6 & [B] \\
\addlinespace
\texttt{N5+\_C2v\_experimental\_Xray} & N$_{5}^{+}$ & C2v & 195{,}8 & [A8] \\
\texttt{N5+\_Vchain} & N$_{5}^{+}$ & C2v & 195{,}9 & [B] \\
\addlinespace
\texttt{N6+\_C2h\_planar} & N$_{6}^{+}$ & C2h & 196{,}0 & [B] \\
\texttt{N6+\_C2v} & N$_{6}^{+}$ & C2v & 196{,}0 & [B] \\
\texttt{N6+\_C2h\_nonplanar} & N$_{6}^{+}$ & C2h & 196{,}0 & [B] \\
\addlinespace
\texttt{N7+\_chain} & N$_{7}^{+}$ & C2v & 168{,}1 & [B] \\
\texttt{N7+\_1\_C2v\_openchain} & N$_{7}^{+}$ & C2v & 168{,}1 & [A11] \\
\addlinespace
\texttt{N9+\_chain} & N$_{9}^{+}$ & C2v & 118{,}3 & [B] \\
\texttt{N9+\_fused\_rings} & N$_{9}^{+}$ & C2v & 156{,}5 & [B] \\
\addlinespace
\texttt{N10+\_chain} & N$_{10}^{+}$ & C1 & 114{,}8 & [B] \\
\texttt{N10+\_rings} & N$_{10}^{+}$ & C1 & 143{,}5 & [B] \\
\end{longtable}


\subsection{Anions}
\begin{longtable}{@{}p{4.6cm}lp{1.4cm}rc@{}}
\caption{Enthalpie de r\'eaction quasi-isodesmique (GFN2-xTB) des compos\'es \textbf{anioniques}, r\'ef\'erenc\'ee \`a NH$_2^-$ ($\Delta H_f=27{,}0$ kcal/mol, valeur exp\'erimentale).}
\label{tab:isodesmic-anion}\\
\toprule
\textbf{Nom} & \textbf{Formule} & \textbf{Sym.} & \textbf{$\Delta H_f$ (kcal/mol)} & \textbf{R\'ef.} \\
\midrule
\endfirsthead
\multicolumn{5}{c}{\small (suite)}\\
\toprule
\textbf{Nom} & \textbf{Formule} & \textbf{Sym.} & \textbf{$\Delta H_f$ (kcal/mol)} & \textbf{R\'ef.} \\
\midrule
\endhead
\bottomrule
\endfoot
\bottomrule
\endlastfoot
\addlinespace
\texttt{N3-\_linear} & N$_{3}^{-}$ & Dinfh & -43{,}4 & [B] \\
\texttt{N3-\_bent} & N$_{3}^{-}$ & C2v & 28{,}8 & [B] \\
\texttt{N3-\_triangle} & N$_{3}^{-}$ & D3h & 47{,}0 & [B] \\
\addlinespace
\texttt{N4-\_rectangle} & N$_{4}^{-}$ & D2h & 5{,}0 & [B] \\
\texttt{N4-\_chain\_planar} & N$_{4}^{-}$ & C2h & 33{,}1 & [B] \\
\addlinespace
\texttt{N5-\_pentagon} & N$_{5}^{-}$ & D5h & -106{,}4 & [B] \\
\texttt{N5-\_pentagon\_D5h\_crosscheck} & N$_{5}^{-}$ & D5h & -106{,}4 & [A2] \\
\texttt{N5-\_D5h\_pentazolate} & N$_{5}^{-}$ & D5h & -106{,}4 & [A7] \\
\texttt{N5-\_bipyramid} & N$_{5}^{-}$ & D3h & 37{,}0 & [B] \\
\addlinespace
\texttt{N6-\_Cs\_chain} & N$_{6}^{-}$ & Cs & -78{,}8 & [B] \\
\texttt{N6-\_C2\_linked\_triangles} & N$_{6}^{-}$ & C2 & 26{,}2 & [B] \\
\addlinespace
\texttt{N7-\_chain} & N$_{7}^{-}$ & C2v & -104{,}3 & [B] \\
\texttt{N7-\_6\_Cs\_6ring\_boat} & N$_{7}^{-}$ & Cs & -95{,}2 & [A11] \\
\addlinespace
\texttt{N8-\_C2h\_chain} & N$_{8}^{-}$ & C2h & -115{,}9 & [B] \\
\texttt{N8-\_C2h\_ZEZ} & N$_{8}^{-}$ & C2h & -115{,}9 & [B] \\
\addlinespace
\texttt{N9-\_chain\_twistedW} & N$_{9}^{-}$ & C2 & -141{,}8 & [B] \\
\texttt{N9-\_chain} & N$_{9}^{-}$ & C2v & -141{,}8 & [B] \\
\texttt{N9-\_ring\_chain} & N$_{9}^{-}$ & Cs & -139{,}1 & [B] \\
\addlinespace
\texttt{N10-\_D2h\_perp\_rings} & N$_{10}^{-}$ & D2h & -186{,}0 & [B] \\
\end{longtable}


\section{R\'ef\'erences}

\begin{description}[leftmargin=1.6cm,itemsep=3pt,style=nextline]
\item[[GS]] Glukhovtsev, Jiao, von Rag\'e Schleyer, \emph{Inorg. Chem.} \textbf{1996}, 35, 7124--7133.
\item[[B]] Fau, Mobita, Wilson, Perera, Bartlett, \emph{Quantum Theory Project report}, University of Florida.
\item[[A1]] Bennett F.R., Maier J.P., Chambaud G., Rosmus P., \emph{Photodissociation, charge and atom transfer processes in electronically excited states of N3+}, Chemical Physics 209 (1996) 275.
\item[[A2]] Fau S., Wilson K.J., Bartlett R.J., \emph{On the Stability of N5+N5-}, J. Phys. Chem. A 2002, 106, 4639.
\item[[A3]] Leininger M.L., Van Huis T.J., Schaefer H.F. III, \emph{Protonated High Energy Density Materials: N4 Tetrahedron and N8 Octahedron}, J. Phys. Chem. A 1997, 101, 4460.
\item[[A4]] Bittererova M., Brinck T., \emph{Theoretical Study of the Triplet N4 Potential Energy Surface}, J. Phys. Chem. A 2000, 104, 11999.
\item[[A5]] Gimarc B.M., Zhao M., \emph{Strain Energies in Homoatomic Nitrogen Clusters N4, N6, and N8}, Inorg. Chem. 1996, 35, 3289.
\item[[A6]] Glukhovtsev M.N., Laiter S., \emph{Thermochemistry of Tetrazete and Tetraazatetrahedrane: A High-Level Computational Study}, J. Phys. Chem. 1996, 100, 1569.
\item[[A7]] Nguyen M.T., Ha T.K., \emph{Decomposition mechanism of the polynitrogen N5 and N6 clusters and their ions}, Chem. Phys. Lett. 335 (2001) 311.
\item[[A8]] Vij A., Wilson W.W., Vij V., Tham F.S., Sheehy J.A., Christe K.O., \emph{Polynitrogen Chemistry. Synthesis, Characterization, and Crystal Structure of Surprisingly Stable Fluoroantimonate Salts of N5+}, J. Am. Chem. Soc. 2001, 123, 6308.
\item[[A9]] Gagliardi L., Evangelisti S., Barone V., Roos B.O., \emph{On the dissociation of N6 into 3 N2 molecules}, Chem. Phys. Lett. 320 (2000) 518.
\item[[A10]] Li Q.S., Hu X.G., Xu W.G., \emph{Structure and stability of N7 cluster}, Chem. Phys. Lett. 287 (1998) 94.
\item[[A11]] Liu Y.D., Zhao J.F., Li Q.S., \emph{Structures and stability of N7+ and N7- clusters}, Theor. Chem. Acc. 107 (2002) 140.
\item[[A12]] Ha T.K., Suleimenov O., Nguyen M.T., \emph{A quantum chemical study of three isomers of N20}, Chem. Phys. Lett. 315 (1999) 327.
\item[[A13]] Li Q.S., Wang L.J., \emph{Theoretical Studies on the Potential Energy Surfaces of N8 Clusters}, J. Phys. Chem. A (recv. Aug 2000).
\item[[A14]] Chung G., Schmidt M.W., Gordon M.S., \emph{An Ab Initio Study of Potential Energy Surfaces for N8 Isomers}, J. Phys. Chem. A 2000, 104, 5647.
\item[[A15]] Gagliardi L., Evangelisti S., Roos B.O., Widmark P.O., \emph{A theoretical study of ten N8 isomers}, J. Mol. Struct. (Theochem) 428 (1998) 1-8.
\item[[A16]] Wang L.J., Warburton P., Mezey P.G., \emph{Theoretical Prediction on the Synthesis Reaction Pathway of N6 (C2h)}, J. Phys. Chem. A 2002, 106, 2748.
\item[[A17]] Wang X., Ren Y., Shuai M.B., Wong N.B., Li W.K., Tian A.M., \emph{Structure and stability of new N7 isomers}, J. Mol. Struct. (Theochem) 538 (2001) 145.
\item[[A18]] Thompson M.D., Bledson T.M., Strout D.L., \emph{Dissociation Barriers for Odd-Numbered Acyclic Nitrogen Molecules N9 and N11}, J. Phys. Chem. A 2002.
\item[[A19]] Gu J.D., Chen K.X., Jiang H.L., Chen J.Z., Ji R.Y., Ren Y., Tian A.M., \emph{N18: a computational investigation}, J. Mol. Struct. (Theochem) 428 (1998) 183.
\item[[A20]] Klapotke T.M., Harcourt R.D., \emph{The interconversion of N12 to N8 and two equivalents of N2}, J. Mol. Struct. (Theochem) 541 (2001) 237.
\end{description}


\end{document}
